\thispagestyle{empty}
\begin{center}
  {\sffamily\small\color{docacento} PROPUESTA TÉCNICA · PLATAFORMA DE AGENTES DE IA}\\[10pt]
  {\sffamily\bfseries\LARGE\color{docink} Agentes como configuración gobernada}\\[6pt]
  {\sffamily\Large\color{docink} Una capa común para crear, operar y controlar los agentes de IA de una comercializadora de energía}\\[14pt]
  {\sffamily\normalsize Daniel Felipe Diab Gonzalez}\\[2pt]
  {\sffamily\small\color{docgris} Octubre de 2026}
\end{center}
\vspace{4pt}

\begin{tesis}
\begin{itemize}[leftmargin=1.1em]
  \item \textbf{El problema.} La empresa ya usa IA en facturación, cobranzas y ventas, pero cada área
  la armó por su cuenta. Nadie sabe cuántos agentes hay, con qué modelo corren, qué datos de
  clientes tocan ni cuánto cuestan.
  \item \textbf{La propuesta.} Una plataforma común donde cada agente es un
  \emph{manifiesto}: una ficha versionada que dice qué hace, de quién es, qué herramientas y
  datos usa y cuánto puede gastar. Sin código y sin prompts escritos a mano.
  \item \textbf{El control.} Un solo camino a producción con tres puntos obligatorios:
  registro, evaluación y gateway. El camino fácil es el camino controlado.
  \item \textbf{La experiencia.} Cualquier área arma su agente conversando con un asistente,
  lo prueba y lo publica; sus usuarios reciben respuestas fluidas, con datos reales.
  \item \textbf{Simple a propósito.} Un runtime (LangGraph), un gateway, una base de datos y una
  API para todos los agentes: menos costo, menos piezas que fallen, menos legado.
  \item \textbf{El plan.} Visibilidad total en 2 semanas, primer agente certificado
  (facturación) en 2 meses y autoservicio para todas las áreas en 4. Desde ahí, cambiar de
  modelo deja de ser un riesgo: cada cambio pasa por evaluación y se revierte en un paso.
\end{itemize}
\end{tesis}

\begin{figure}[h]
  \centering
  \begin{tikzpicture}[x=1cm,y=1cm, cp/.style={ctrl, text width=2.6cm, minimum height=1.25cm}]
\node[neut, text width=2.0cm, minimum height=1.25cm] (idea) at (0,0) {\textbf{Idea de un área}\\``quiero un agente que\ldots''};
\node[cp] (c1) at (3.55,0) {\textbf{1. Registro}\\¿qué es, de quién es y qué toca?};
\node[cp] (c2) at (7.05,0) {\textbf{2. Compuerta}\\¿funciona y es seguro para su riesgo?};
\node[exec, text width=2.6cm, minimum height=1.25cm] (c3) at (10.55,0) {\textbf{3. Gateway}\\¿con qué modelo, qué datos y cuánto cuesta?};
\node[exec, fill=execbd!30, text width=2.0cm, minimum height=1.25cm] (prod) at (14.0,0) {\textbf{Producción}\\clientes y equipos};
\draw[flecha, line width=1.1pt, draw=docacento] (idea) -- (c1);
\draw[flecha, line width=1.1pt, draw=docacento] (c1) -- (c2);
\draw[flecha, line width=1.1pt, draw=docacento] (c2) -- (c3);
\draw[flecha, line width=1.1pt, draw=docacento] (c3) -- (prod);
\foreach \c/\t in {c1/llena el inventario\\de agentes, c2/deja evidencia\\de calidad, c3/mide costo y uso\\de datos}{
  \node[plain, text=docacento, font=\sffamily\tiny\itshape, anchor=north] at ($(\c.south)+(0,-0.08)$) {\t};}
\draw[decorate, decoration={brace, amplitude=5pt}, draw=docgris] ($(c1.north west)+(0,0.12)$) -- node[plain, above=5pt, text=docgris]{camino pavimentado: tres puntos de control obligatorios} ($(c3.north east)+(0,0.12)$);
\node[risk, text width=5.2cm] (atajo) at (7.05,-2.35) {\textbf{Atajo}: herramienta externa con llaves propias\\sin llaves corporativas · sin herramientas internas · sin datos de clientes};
\draw[flechad, draw=riskbd] (idea.south) |- (atajo.west);
\node[circle, draw=riskbd, fill=white, text=riskbd, inner sep=0.5pt, font=\sffamily\bfseries\small] at ($(atajo.east)+(0.35,0)$) {$\times$};
\end{tikzpicture}

  \caption{La propuesta en un diagrama: el camino más fácil a producción es también el único
  que tiene llaves, herramientas internas y datos de clientes.}
  \label{fig:tesis}
\end{figure}



\newpage
{\small\setlength{\parskip}{1pt}\tableofcontents}
\vspace{6pt}
